% bibliography of Volume 3: the 235 of 659 entries of ch/bib.tex it cites; generated
\begin{thebibliography}{999}\small\setlength{\itemsep}{0pt}\setlength{\parsep}{0pt}\setlength{\parskip}{0pt}
% verified: CrossRef 10.1002/nme.995 2026-10-09

% verified: arXiv 2608.13637 2026-10-11 (version 2 of 19 August 2026; title, authors and the constants 2/3 and 0.6725 checked in the abstract)
\bibitem{AF26} L. Alp\"oge and R. Furman. More than two thirds of the zeta zeros are simple and on the critical line. arXiv:2608.13637, 2026.
% verified: CrossRef 10.2996/kmj/1490083225 2026-11-02 (fields; the theorem itself not read)
\bibitem{Akatsuka17} H. Akatsuka. The Euler product for the Riemann zeta-function in the critical strip. \emph{Kodai Mathematical Journal} 40 (2017), 79--101.
% verified: zbMATH 1097.11041 2026-09-26
\bibitem{BaezDuarte03} L. B\'aez-Duarte. A strengthening of the Nyman--Beurling criterion for the Riemann hypothesis. \emph{Atti della Accademia Nazionale dei Lincei, Rendiconti Lincei, Matematica e Applicazioni} 14 (2003), 5--11.
% verified: CrossRef 10.1155/2010/714147 2026-10-26 (article number, year; arXiv math/0307214); the abstract says only that the paper gives a rigorous proof of Maslanka's representation, so no decay or convergence rate is credited to it
\bibitem{BaezDuarte10} L. B\'aez-Duarte. On Maslanka's representation for the Riemann zeta function. \emph{International Journal of Mathematics and Mathematical Sciences} (2010), Article ID 714147. doi:10.1155/2010/714147, arXiv:math/0307214.
% verified: arXiv 1306.0856 abstract page 2026-10-25 (journal, volume, issue, pages and DOI as recorded there)
\bibitem{BLM14} H. M. Bui, S. J. Lester and M. B. Milinovich. On Balazard, Saias, and Yor's equivalence to the Riemann Hypothesis. \emph{Journal of Mathematical Analysis and Applications} 409 (2014), no. 1, 244--253. doi:10.1016/j.jmaa.2013.06.063, arXiv:1306.0856.
% verified: CrossRef 10.1090/s0025-5718-07-01999-0 2026-10-11 (Math. Comp. 76, no. 260, 2045-2049); abstract (CrossRef and OpenAlex): zeros off the critical line of the Dirichlet series of Davenport and Heilbronn
\bibitem{BSO07} E. Balanzario and J. S\'anchez-Ortiz. Zeros of the Davenport--Heilbronn counterexample. \emph{Mathematics of Computation} 76 (2007), 2045--2049.
\bibitem{BNKS82} O. E. Barndorff-Nielsen, J. Kent and M. S{\o}rensen. Normal variance-mean mixtures and $z$ distributions. \emph{International Statistical Review} 50 (1982), 145 ff.
% verified: zbMATH 1102.60013 2026-10-10 (not in CrossRef; the DOI suggested earlier, 10.3150/bj/1137421635, belongs to another Bernoulli paper)
\bibitem{BNMS06} O. E. Barndorff-Nielsen, M. Maejima and K. Sato. Some classes of multivariate infinitely divisible distributions admitting stochastic integral representations. \emph{Bernoulli} 12 (2006), 1--33.
% verified: zbMATH 1024.60022 2026-10-10 (not in CrossRef)
\bibitem{BNT02} O. E. Barndorff-Nielsen and S. Thorbj{\o}rnsen. Self-decomposability and L\'evy processes in free probability. \emph{Bernoulli} 8 (2002), 323--366.
% verified: CrossRef 10.1090/mcom/3599 2026-10-10 (title, journal, volume 90, pages 1455 to 1482, issued 2021, authors Bennett, Martin, O'Bryant, Rechnitzer); arXiv 2005.02989 abstract checked for the bound used in Theorem thm:dbeta:minall(v): for conductor q > 1, T >= 5/7 and l = log(q(T+2)/(2 pi)) > 1.567, |N(T,chi) - ((T/pi) log(qT/(2 pi e)) - chi(-1)/4)| <= 0.22737 l + 2 log(1+l) - 0.5
\bibitem{BMOR21} M. A. Bennett, G. Martin, K. O'Bryant and A. Rechnitzer. Counting zeros of Dirichlet $L$-functions. \emph{Mathematics of Computation} 90 (2021), 1455--1482.
% verified: CrossRef 10.2307/121080 2026-10-10 (CrossRef gives the first page only; pages 1023 to 1060 from zbMATH 0945.46046; CrossRef lists P. Biane, the author of the appendix, as a third author)
\bibitem{BP99} H. Bercovici and V. Pata. Stable laws and domains of attraction in free probability theory (with an appendix by P. Biane). \emph{Annals of Mathematics} 149 (1999), 1023--1060.
% verified: CrossRef 10.1006/aima.1998.1797 2026-10-25
\bibitem{BSY99} M. Balazard, E. Saias and M. Yor. Notes sur la fonction $\zeta$ de Riemann, 2. \emph{Advances in Mathematics} 143 (1999), 284--287.
% verified: arXiv 2606.24405v1 (math-ph, 23 June 2026), title, authors and the statements of Proposition 1, Lemma 2 and Theorem 3 checked against the text of the preprint; unrefereed
\bibitem{BagarelloKuzel26} F. Bagarello and S. Ku\.{z}el. On the Berry--Keating operator. arXiv:2606.24405, preprint, 2026.
% verified: CrossRef 10.1137/s0036144598347497 2026-10-09
\bibitem{BerryKeating99} M. V. Berry and J. P. Keating. The Riemann zeros and eigenvalue asymptotics. \emph{SIAM Review} 41 (1999), 236--266.
% verified: CrossRef 10.1073/pnas.41.5.312 2026-09-26
\bibitem{Beurling55} A. Beurling. A closure problem related to the Riemann zeta-function. \emph{Proceedings of the National Academy of Sciences} 41 (1955), 312--314.
% verified: CrossRef 10.1090/s0273-0979-01-00912-0 2026-09-26
\bibitem{BPY01} P. Biane, J. Pitman and M. Yor. Probability laws related to the Jacobi theta and Riemann zeta functions, and Brownian excursions. \emph{Bulletin of the American Mathematical Society} 38 (2001), 435--465.
% verified: CrossRef 10.1017/cbo9780511721434 2026-10-05 (Theorem 4.12.9 is de Bruijn's exponential Tauberian theorem, as cited by Aurzada and Dereich, Ann. Probab. 37 (2009), arXiv:0805.1330, reference [4]; the 1989 paperback keeps the numbering)
\bibitem{BGT87} N. H. Bingham, C. M. Goldie and J. L. Teugels. \emph{Regular Variation}. Encyclopedia of Mathematics and its Applications 27, Cambridge University Press, 1987.
% verified: CrossRef 10.1016/j.aim.2020.107240, arXiv 2004.11501 2026-11-02
\bibitem{BDV20} F. Broucke, G. Debruyne and J. Vindas. Beurling integers with RH and large oscillation. \emph{Advances in Mathematics} 370 (2020), 107240.
% verified: zbMATH 45.0716.02 2026-09-26
\bibitem{BL14} H. Bohr and E. Landau. Sur les z\'eros de la fonction $\zeta(s)$ de Riemann. \emph{Comptes Rendus de l'Acad\'emie des Sciences} 158 (1914), 106--110.
% verified: CrossRef 10.1006/jnth.1999.2392 2026-09-26
\bibitem{BL99} E. Bombieri and J. C. Lagarias. Complements to Li's criterion for the Riemann hypothesis. \emph{Journal of Number Theory} 77 (1999), 274--287.
% verified: zbMATH 1008.11034 2026-09-26
\bibitem{BombieriStepanov} E. Bombieri. Counting points on curves over finite fields (d'apr\`es S. A. Stepanov). In \emph{S\'eminaire Bourbaki}, expos\'e 430 (1972/73), Lecture Notes in Mathematics 383, Springer, 1974, 234--241. % verified: CrossRef 10.1007/bfb0057311 and zbMATH (Exp. No. 430, LNM 383, 1974) 2026-11-10
\bibitem{Bombieri} E. Bombieri. Remarks on Weil's quadratic functional in the theory of prime numbers, I. \emph{Rendiconti Lincei. Matematica e Applicazioni} (9) 11 (2000), 183--233.
% verified: CrossRef 10.1007/978-1-4612-2948-3 2026-09-26
\bibitem{Bondesson} L. Bondesson. \emph{Generalized Gamma Convolutions and Related Classes of Distributions and Densities}. Lecture Notes in Statistics 76, Springer, 1992.
% verified: CrossRef 10.1111/1467-9965.00035 2026-10-09
\bibitem{BGK97} M. Broadie, P. Glasserman and S. Kou. A continuity correction for discrete barrier options. \emph{Mathematical Finance} 7 (1997), 325--349.
% New keys of editor E27 (ch/c27.tex, cg1 to cg3). Fields are those of the source bibliography, nothing added.
% verified: pending
\bibitem{ComtetMajumdar05} A. Comtet and S. N. Majumdar. Precise asymptotics for a random walker's maximum. \emph{Journal of Statistical Mechanics} (2005), P06013; arXiv cond-mat/0506195.
% verified: zbMATH 0966.11039 2026-09-26
\bibitem{CL00} J. B. Conrey and X.-J. Li. A note on some positivity conditions related to zeta and $L$-functions. \emph{International Mathematics Research Notices} 2000 (2000), 929--940.
% verified: CrossRef 10.1007/bf01482072 2026-09-26
\bibitem{Cramer} H. Cram\'er. Ein Mittelwertsatz in der Primzahltheorie. \emph{Mathematische Zeitschrift} 12 (1922), 147--153.
% verified: CrossRef 10.1090/s0002-9947-1986-0846596-4 2026-09-27
\bibitem{CNV86} G. Csordas, T. S. Norfolk and R. S. Varga. The Riemann hypothesis and the Tur\'an inequalities. \emph{Transactions of the American Mathematical Society} 296 (1986), 521--541.
\bibitem{CSV94} G. Csordas, W. Smith and R. S. Varga. Lehmer pairs of zeros, the de Bruijn-Newman constant $\Lambda$, and the Riemann hypothesis. \emph{Constructive Approximation} 10 (1994), 107--129.
% verified: CrossRef 10.1007/bf02788653 2026-10-23
\bibitem{CurrySchoenberg66} H. B. Curry and I. J. Schoenberg. On P\'olya frequency functions IV: the fundamental spline functions and their limits. \emph{Journal d'Analyse Math\'ematique} 17 (1966), 71--107.
% verified: zbMATH 1002.11001 2026-09-26
\bibitem{Davenport} H. Davenport. \emph{Multiplicative Number Theory}, third edition, revised by H. L. Montgomery. Graduate Texts in Mathematics 74, Springer, 2000.
\bibitem{Devroye09} L. Devroye. On exact simulation algorithms for some distributions related to Jacobi theta functions. \emph{Statistics and Probability Letters} 79 (2009), 2251--2259.
% verified: CrossRef 10.1112/jlms/s1-11.3.181 (part I, s1-11 no. 3, 181-185, July 1936) and 10.1112/jlms/s1-11.4.307 (part II, s1-11 no. 4, 307-312, October 1936), 2026-10-11; text not seen, so the part that contains the example of period 5 is unchecked and both are cited
\bibitem{DH36} H. Davenport and H. Heilbronn. On the zeros of certain Dirichlet series I, II. \emph{Journal of the London Mathematical Society} 11 (1936), 181--185 and 307--312.
% verified: CrossRef 10.1215/s0012-7094-50-01720-0 2026-09-26; pages 197 to 226 from zbMATH 0038.23302, 2026-10-02
\bibitem{deBruijn} N. G. de Bruijn. The roots of trigonometric integrals. \emph{Duke Mathematical Journal} 17 (1950), 197--226.
% verified: zbMATH 0327.60029 2026-09-26
\bibitem{DymMcKean76} H. Dym and H. P. McKean. \emph{Gaussian Processes, Function Theory, and the Inverse Spectral Problem}. Probability and Mathematical Statistics 31, Academic Press, 1976.
% verified: CrossRef 10.1063/1.1703862 2026-10-09
\bibitem{Dyson62} F. J. Dyson. A Brownian-motion model for the eigenvalues of a random matrix. \emph{Journal of Mathematical Physics} 3 (1962), 1191--1198.
% verified: pending
\bibitem{Edrei53} A. Edrei. On the generating functions of totally positive sequences II. \emph{Journal d'Analyse Math\'ematique} 2 (1952), 104--109.
% verified: zbMATH 0219.60003 2026-09-26
\bibitem{Feller71} W. Feller. \emph{An Introduction to Probability Theory and Its Applications}, volume II, second edition. Wiley, 1971.
% verified: CrossRef 10.1073/pnas.1902572116 2026-10-01
\bibitem{GORZ19} M. Griffin, K. Ono, L. Rolen and D. Zagier. Jensen polynomials for the Riemann zeta function and other sequences. \emph{Proceedings of the National Academy of Sciences} 116 (2019), 11103--11110.
% verified: CrossRef 10.1007/bf02465533 2026-10-12 (English translation of Liet. Mat. Rink. 39 (1999), 40 to 51; the publisher's page returned a client challenge, so the paper itself was not read)
\bibitem{Grigelionis99} B. Grigelionis. Processes of Meixner type. \emph{Lithuanian Mathematical Journal} 39 (1999), 33--41.
% verified: CrossRef 10.1515/crll.1914.144.114 2026-10-11; text not checked (cited through BS18, Theorem 10 (i), reference [8])
\bibitem{Grommer14} J. Grommer. Ganze transzendente Funktionen mit lauter reellen Nullstellen. \emph{Journal f\"ur die reine und angewandte Mathematik} 144 (1914), 114--166.
% verified: pending
\bibitem{Guinand48} A. P. Guinand. A summation formula in the theory of prime numbers. \emph{Proceedings of the London Mathematical Society} 50 (1948), 107--119.
% verified: CrossRef 10.4007/annals.2026.203.2.6 2026-09-27 (pages 623 to 675 from the journal; CrossRef gives no pages)
\bibitem{GM26} L. Guth and J. Maynard. New large value estimates for Dirichlet polynomials. \emph{Annals of Mathematics} 203 (2026), 623--675.
% verified: arXiv 0910.5228 2026-10-16 (title and authors checked on the abstract page)
\bibitem{Haglund09} J. Haglund. Some conjectures on the zeros of approximates to the Riemann $\Xi$-function and incomplete gamma functions. arXiv:0910.5228, 2009.
% verified: zbMATH 46.0498.01 2026-09-26
\bibitem{HL16} G. H. Hardy and J. E. Littlewood. Contributions to the theory of the Riemann zeta-function and the theory of the distribution of primes. \emph{Acta Mathematica} 41 (1916), 119--196.
% verified: CrossRef 10.1016/j.jnt.2021.06.032 2026-10-11 (title, journal, volume 235, pages 219 to 241, issued June 2022, authors Hasanalizade, Shen, Wong); arXiv 2107.06506v1 read for (1.7), (2.7), Corollaries 1.2 and 1.4 and footnote 2 (the error in Trudgian 2014)
\bibitem{HSW22} E. Hasanalizade, Q. Shen and P.-J. Wong. Counting zeros of the Riemann zeta function. \emph{Journal of Number Theory} 235 (2022), 219--241.
% verified: CrossRef 10.1007/bf01565437 2026-10-09
\bibitem{Hecke36} E. Hecke. \"Uber die Bestimmung Dirichletscher Reihen durch ihre Funktionalgleichung. \emph{Mathematische Annalen} 112 (1936), 664--699.
% verified: CrossRef 10.1137/1101021 2026-10-23 (not in the source; supplied from memory on the request of review-B m13: Ibragimov, On the composition of unimodal distributions, Theory of Probability and its Applications 1 (1956), 255--260; check before use)
\bibitem{Ibragimov56} I. A. Ibragimov. On the composition of unimodal distributions. \emph{Theory of Probability and its Applications} 1 (1956), 255--260.
% verified: CrossRef 10.2977/PRIMS/1216238306 2026-11-08
\bibitem{Ihara08} Y. Ihara. On $M$-functions closely related to the distribution of $L'/L$-values. \emph{Publications of the RIMS} 44 (2008), 893--954.
% E3 (ht5), keys used in ch/ht5.tex; entries below not yet in ch/bib.tex
% verified: pending
\bibitem{Karlin68} S. Karlin. \emph{Total Positivity}, Volume I. Stanford University Press, 1968.
% UNCHECKED: not found in CrossRef or OpenAlex
\bibitem{KershnerWintner37} R. Kershner and A. Wintner. On the asymptotic distribution of $\zeta'/\zeta(s)$ in the critical strip. \emph{American Journal of Mathematics} 59 (1937), 673--678.
% verified: CrossRef 10.1002/mana.19770770116 2026-11-10
\bibitem{KL77} M. G. Krein and H. Langer. \"Uber einige Fortsetzungsprobleme, die eng mit der Theorie hermitescher Operatoren im Raume $\Pi_\kappa$ zusammenh\"angen. I. Einige Funktionenklassen und ihre Darstellungen. \emph{Mathematische Nachrichten} 77 (1977), 187--236.
% verified: pending (the source gives year 1981 and Lecture Notes in Math. 923, pages 235 to 259; the reader A1 suggests that the volume appeared in 1982, to be confirmed before the entry is used)
\bibitem{KotaniWatanabe82} S. Kotani and S. Watanabe. Krein's spectral theory of strings and generalized diffusion processes. \emph{Lecture Notes in Mathematics} 923 (1982), 235--259.
% verified: CrossRef 10.1086/111605 2026-10-24
\bibitem{Lucy74} L. B. Lucy. An iterative technique for the rectification of observed distributions. \emph{Astronomical Journal} 79 (1974), 745--754.
% verified: arXiv math-ph/0105007 abstract page 2026-10-26 (submitted 4 May 2001; no journal version found)
\bibitem{Maslanka01} K. Ma\'slanka. Hypergeometric-like representation of the zeta-function of Riemann. arXiv:math-ph/0105007, 2001.
% verified: CrossRef 10.4064/aa166-2-4 2026-10-25
\bibitem{MSZ14} Yu. Matiyasevich, F. Saidak and P. Zvengrowski. Horizontal monotonicity of the modulus of the zeta function, Dirichlet $L$-functions, and related functions. \emph{Acta Arithmetica} 166 (2014), 189--200.
% verified: pending (not found: CrossRef and OpenAlex searched by title)
\bibitem{NCKolSimon05} B. Simon. \emph{Functional Integration and Quantum Physics}. Second edition, AMS Chelsea Publishing, 2005. % verified: Open Library ISBN 0821835823 (AMS Chelsea Pub.) 2026-11-09
\bibitem{NCKolWuSprung93} H. Wu and D. W. L. Sprung. Riemann zeros and a fractal potential. \emph{Physical Review E} 48 (1993), 2595--2598. % verified: DOI 10.1103/physreve.48.2595 2026-11-09
% verified: pending
\bibitem{P26cg} N. G. Polson. The Comtet and Majumdar identity, the $2\sin^2$ lemma, and approximating families for the reciprocal Xi function. Manuscript, October 2026.
% verified: pending (the manuscript is the source of this chapter, untracked, see misc/reviews/nick-lvj_sech/source5)
\bibitem{P26es} N. G. Polson. The explicit formula as a deconvolution. Manuscript, October 2026.
% verified: pending (title and sole author from the source file; defined here only if the editor of ch/c25.tex has not added it)
\bibitem{P26evt} N. G. Polson. Extreme values, hitting times and a path to RH. Manuscript, October 2026.
% verified: pending (source note of the second half of Chapter 28; the note is headed "Progress on the Open Problems", a note for the book's editor; no journal)
\bibitem{P26hermite} N. G. Polson. Progress on the open problems. Manuscript, October 2026.
% verified: pending
\bibitem{P26ht} N. G. Polson. Hitting times and the reciprocal Xi function. Manuscript, October 2026.
% verified: pending
\bibitem{P26lo} N. G. Polson. The Thorin measure at level one and the backward Cauchy flow. Manuscript, October 2026.
% verified: pending (source note of the first half of Chapter 28; title and date from the note, no journal)
\bibitem{P26trunc} N. G. Polson. Thorin truncation and the zeros of $L(s,\chi_4)$. Manuscript, October 2026.
% verified: pending
\bibitem{P26wt} N. G. Polson. Weil positivity as the Thorin condition. Manuscript, October 2026.
% verified: CrossRef 10.1214/18-EJS1430 2026-11-08
\bibitem{PadillaPolsonScott18} O. H. Madrid Padilla, N. G. Polson and J. G. Scott. A deconvolution path for mixtures. \emph{Electronic Journal of Statistics} 12 (2018), 1717--1751.
% verified: pending
\bibitem{PolsonEuler26} N. G. Polson. The Euler product of $1/\zeta$ as a L\'evy measure, the $2\sin^2$ lemma, and the reciprocal Xi function as a generalized gamma convolution. Note, October 2026.
% verified: pending
\bibitem{PolsonResidue26} N. G. Polson. A residue expansion of the reciprocal Xi function and generalized gamma convolution methods. Note, October 2026.
% verified: the note itself, title from its first page (misc/reviews/nick-rxe/source, N:41) 2026-10-26
\bibitem{PolsonRXi26} N. G. Polson. The reciprocal Xi function as the Laplace transform of an EGGC: mixing laws, the Euler product, Selberg and Barnes. Note, October 2026.
% verified: CrossRef 10.1111/rssb.12130 2026-11-08
\bibitem{PolsonScott16} N. G. Polson and J. G. Scott. Mixtures, envelopes and hierarchical duality. \emph{Journal of the Royal Statistical Society, Series B} 78 (2016), 701--727.
% verified: CrossRef 10.1080/01621459.2013.829001 2026-10-22 (J. Amer. Statist. Assoc. 108 (2013), 1339-1349 confirmed; authors Polson, Scott, Windle)
\bibitem{PSW13} N. G. Polson, J. G. Scott and J. Windle. Bayesian Inference for Logistic Models Using P\'olya--Gamma Latent Variables. \emph{Journal of the American Statistical Association} 108 (2013), 1339--1349.
% verified: CrossRef 10.2307/2372402 2026-11-08 (start page 97; end page not in the record)
\bibitem{Putnam54} C. R. Putnam. On the non periodicity of the zeros of the Riemann zeta function. \emph{American Journal of Mathematics} 76 (1954), 97.
% verified: CrossRef 10.1364/JOSA.62.000055 2026-10-24
\bibitem{Richardson72} W. H. Richardson. Bayesian based iterative method of image restoration. \emph{Journal of the Optical Society of America} 62 (1972), 55--59.
% verified: CrossRef 10.1109/tmi.1982.4307558 2026-10-24
\bibitem{ShVa82} L. A. Shepp and Y. Vardi. Maximum likelihood reconstruction for emission tomography. \emph{IEEE Transactions on Medical Imaging} 1 (1982), 113--122.
% verified: pending
\bibitem{Spitzer56} F. Spitzer. A combinatorial lemma and its application to probability theory. \emph{Transactions of the American Mathematical Society} 82 (1956), 323--339.
% verified: zbMATH 0241.10026 2026-09-27 (year 1972 as in zbMATH and MR 292774; CrossRef 10.1007/bf01418933 dates the issue June 1971, before the paper's receipt on 4 November 1971)
\bibitem{Huxley72} M. N. Huxley. On the difference between consecutive primes. \emph{Inventiones Mathematicae} 15 (1972), 164--170.
% verified: zbMATH 0006.39701 2026-10-17 (author, title, Cambridge Tracts 30, Cambridge University Press, 1932; no CrossRef record found; zbMATH 0715.11045 gives the same series number 30 for the 1990 Cambridge Mathematical Library reprint)
\bibitem{Ingham32} A. E. Ingham. \emph{The Distribution of Prime Numbers}. Cambridge Tracts in Mathematics and Mathematical Physics 30, Cambridge University Press, 1932.
% verified: zbMATH 0025.02704 2026-09-26
\bibitem{Ingham40} A. E. Ingham. On the estimation of $N(\sigma,T)$. \emph{Quarterly Journal of Mathematics} 11 (1940), 291--292.
% verified: CrossRef 10.2307/2371685 2026-10-01; pages 313 to 319 from zbMATH 0063.02974, 2026-10-02
\bibitem{Ingham42} A. E. Ingham. On two conjectures in the theory of numbers. \emph{American Journal of Mathematics} 64 (1942), 313--319.
% verified: zbMATH 1059.11001 2026-09-27
\bibitem{IwaniecKowalski04} H. Iwaniec and E. Kowalski. \emph{Analytic Number Theory}. American Mathematical Society Colloquium Publications 53, American Mathematical Society, 2004.
% verified: CrossRef 10.1109/tc.2017.2690633 2026-09-26
\bibitem{Arb} F. Johansson. Arb: efficient arbitrary-precision midpoint-radius interval arithmetic. \emph{IEEE Transactions on Computers} 66 (2017), 1281--1292.
% verified: CrossRef 10.1090/s0273-0979-99-00766-1 2026-10-17; zbMATH 0921.11047 and OpenAlex agree
\bibitem{KatzSarnak99} N. M. Katz and P. Sarnak. Zeroes of zeta functions and symmetry. \emph{Bulletin of the American Mathematical Society} 36 (1999), 1--26.
\bibitem{Keiper92} J. B. Keiper. Power series expansions of Riemann's $\xi$ function. \emph{Mathematics of Computation} 58 (1992), 765--773.
% verified: CrossRef 10.1007/bf00538895 2026-09-26
\bibitem{Kent80} J. T. Kent. Eigenvalue expansions for diffusion hitting times. \emph{Zeitschrift f\"ur Wahrscheinlichkeitstheorie und Verwandte Gebiete} 52 (1980), 309--319.
% verified: CrossRef 10.1214/aop/1176988867 2026-10-09 (pages 442 to 452 from zbMATH 0796.60017)
\bibitem{Kristiansen94} G. K. Kristiansen. A proof of Steutel's conjecture. \emph{Annals of Probability} 22 (1994), 442--452.
% verified: CrossRef 10.1006/jath.1999.3316 2026-10-10
\bibitem{KVA99} A. B. J. Kuijlaars and W. Van Assche. The asymptotic zero distribution of orthogonal polynomials with varying recurrence coefficients. \emph{Journal of Approximation Theory} 99 (1999), 167--197.
% verified: CrossRef 10.4064/aa-89-3-217-234 2026-09-26
\bibitem{Lagarias} J. C. Lagarias. On a positivity property of the Riemann $\xi$-function. \emph{Acta Arithmetica} 89 (1999), 217--234; correction, \emph{Acta Arithmetica} 116 (2005), 293--294. % correction verified: web search and zbMATH review of Zbl 1026.11070, 2026-11-10
\bibitem{Garunkstis02} R. Garunk\v{s}tis. On a positivity property of the Riemann $\xi$-function. \emph{Lithuanian Mathematical Journal} 42 (2002), 140--145. % verified: zbMATH (Zbl 1026.11070), DOI 10.1023/A:1016158125685 2026-11-10
% verified: CrossRef 10.1111/j.1467-842x.1979.tb01143.x 2026-09-26
\bibitem{LVJ} C. D. Lai and D. Vere-Jones. Odd man out: the Meixner hypergeometric distribution. \emph{Australian Journal of Statistics} 21 (1979), 256--265.
% verified: zbMATH 0152.06703 2026-10-16 (book; series and year checked)
\bibitem{Levin64} B. Ja. Levin. \emph{Distribution of Zeros of Entire Functions}. Translations of Mathematical Monographs 5, American Mathematical Society, 1964.
% verified: CrossRef 10.1006/jnth.1997.2137 2026-09-26
\bibitem{Li97} X.-J. Li. The positivity of a sequence of numbers and the Riemann hypothesis. \emph{Journal of Number Theory} 65 (1997), 325--333.
% verified: CrossRef 10.1090/tran/7249 2026-09-26
\bibitem{LPS18} A. Lindner, L. Pan and K. Sato. On quasi-infinitely divisible distributions. \emph{Transactions of the American Mathematical Society} 370 (2018), 8483--8520.
% verified: CrossRef 10.1017/s0305004100014225 2026-10-16; the record gives the later journal title, Mathematical Proceedings of the Cambridge Philosophical Society
\bibitem{Littlewood24} J. E. Littlewood. On the zeros of the Riemann zeta-function. \emph{Proceedings of the Cambridge Philosophical Society} 22 (1924), 295--318.
\bibitem{Liu26a} B. Liu. Improved zero-free half-planes for the quasi-Riemann hypothesis. arXiv:2610.12234 [math.NT], submitted 8 October 2026, v2 12 October 2026. Preprint, not refereed. % verified: arXiv listing, reference check of Nick Polson of 21 October 2026
\bibitem{Liu26b} B. Liu. Zero-free half-planes for Hecke $L$-functions over number fields. arXiv:2610.26835, submitted 19 October 2026. Preprint, not refereed. % verified: arXiv listing, reference check of Nick Polson of 21 October 2026
% verified: CrossRef 10.1111/j.2517-6161.1993.tb01925.x 2026-10-24 (pages corrected to 569--598)
\bibitem{VardiLee93} Y. Vardi and D. Lee. From image deblurring to optimal investments: maximum likelihood solutions for positive linear inverse problems. \emph{Journal of the Royal Statistical Society, Series B} 55 (1993), 569--598.
% verified: pending
\bibitem{Voros10} A. Voros. \emph{Zeta Functions over Zeros of Zeta Functions}. Lecture Notes of the Unione Matematica Italiana 8, Springer, 2010.
\bibitem{XuSang26} Y. Xu and H. Sang. Amplified moments and a refinement of the quasi-Riemann hypothesis. arXiv:2610.15291, October 2026. Preprint, not refereed. % verified: arXiv listing, reference check of Nick Polson of 21 October 2026
\bibitem{Xylouris11} T. Xylouris. On the least prime in an arithmetic progression and estimates for the zeros of Dirichlet $L$-functions. \emph{Acta Arithmetica} 150 (2011), no.~1, 65--91. doi:10.4064/aa150-1-4. % verified: IMPAN page, reference check of Nick Polson of 21 October 2026; the Linnik constant of this paper is $L=5.2$, the value $5$ is from the later Bonn dissertation
\bibitem{Lukacs68} E. Lukacs. Contributions to a problem of D. van Dantzig. \emph{Theory of Probability and its Applications} 13 (1968), 116--127.
% verified: CrossRef 10.1090/pspum/024/9944 2026-09-27
\bibitem{Montgomery73} H. L. Montgomery. The pair correlation of zeros of the zeta function. In \emph{Analytic Number Theory}, H. G. Diamond, ed., Proceedings of Symposia in Pure Mathematics 24, American Mathematical Society, 1973, 181--193.
% verified: CrossRef 10.3150/13-BEJ581 2026-10-01 (pages 604 to 617 from the journal line of the arXiv reprint 1504.03438; CrossRef gives no pages)
\bibitem{Nakamura15} T. Nakamura. A complete Riemann zeta distribution and the Riemann hypothesis. \emph{Bernoulli} 21 (2015), 604--617.
% verified: CrossRef 10.1007/bf01608542 2026-10-01
\bibitem{Newman75} C. M. Newman. Inequalities for Ising models and field theories which obey the Lee-Yang theorem. \emph{Communications in Mathematical Physics} 41 (1975), 1--9.
% verified: CrossRef 10.1090/s0002-9939-1976-0434982-5 2026-09-25
\bibitem{Newman1976} C. M. Newman. Fourier transforms with only real zeros. \emph{Proceedings of the American Mathematical Society} 61 (1976), 245--251.
% verified: zbMATH 0037.35401 2026-09-26
\bibitem{Nyman50} B. Nyman. \emph{On the One-Dimensional Translation Group and Semi-Group in Certain Function Spaces}. Thesis, University of Uppsala, 1950.
% verified: CrossRef 10.1090/s0025-5718-1987-0866115-0 2026-09-25
\bibitem{Odlyzko87} A. M. Odlyzko. On the distribution of spacings between zeros of the zeta function. \emph{Mathematics of Computation} 48 (1987), 273--308.
% verified: CrossRef 10.1090/conm/290/04578 2026-09-27
\bibitem{Odlyzko01} A. M. Odlyzko. The $10^{22}$-nd zero of the Riemann zeta function. In \emph{Dynamical, Spectral, and Arithmetic Zeta Functions}, M. L. Lapidus and M. van Frankenhuysen, eds., Contemporary Mathematics 290, American Mathematical Society, 2001, 139--144.
% verified: URL https://www-users.cse.umn.edu/~odlyzko/zeta_tables/ 2026-09-27
\bibitem{OdlyzkoTables} A. M. Odlyzko. \emph{Tables of zeros of the Riemann zeta function}. Online, \url{https://www-users.cse.umn.edu/~odlyzko/zeta_tables/}, accessed 27 September 2026.
% verified: publisher URL https://dlmf.nist.gov/help/cite 2026-09-27
\bibitem{DLMF} F. W. J. Olver, A. B. Olde Daalhuis, D. W. Lozier, B. I. Schneider, R. F. Boisvert, C. W. Clark, B. R. Miller, B. V. Saunders, H. S. Cohl and M. A. McClain, eds. \emph{NIST Digital Library of Mathematical Functions}. \url{https://dlmf.nist.gov/}, Release 1.2.8 of 2026-09-15.
\bibitem{OAI26a} OpenAI. The quasi-Riemann hypothesis: a zero-free half-plane $\re(s)>7/8$. Preprint dated 30 September 2026, released 6 October 2026 with a formalisation in Lean, \url{https://github.com/openai/math} (family 003), accessed 7 October 2026.
\bibitem{OAI26b} OpenAI. The quasi-Riemann hypothesis. Preprint dated 5 October 2026, released 6 October 2026, \url{https://github.com/openai/math} (family 003), accessed 7 October 2026.
% verified: zbMATH 1345.11064 2026-09-26
\bibitem{Platt16} D. J. Platt. Numerical computations concerning the GRH. \emph{Mathematics of Computation} 85 (2016), 3009--3027.
% verified: CrossRef 10.1112/blms.12460 2026-09-26
\bibitem{PT21} D. J. Platt and T. S. Trudgian. The Riemann hypothesis is true up to $3\cdot10^{12}$. \emph{Bulletin of the London Mathematical Society} 53 (2021), 792--797.
% verified: arXiv 1708.02653v24 2026-09-27
\bibitem{P17} N. G. Polson. On Hilbert's 8th problem. arXiv:1708.02653v24, 2026.
% verified: arXiv 1708.02653v5 2026-09-27
\bibitem{P17v5} N. G. Polson. On Hilbert's 8th problem. arXiv:1708.02653v5, 2017.
% verified: arXiv 1804.10043 2026-10-14 (title of the current versions v7 to v9; v1 to v6 are titled 'van Dantzig Pairs, Wald Couples and Hadamard Factorisation'; no journal version found in CrossRef)
\bibitem{P18} N. G. Polson. Riemann, Thorin, van Dantzig pairs, Wald couples and Hadamard factorisation. arXiv:1804.10043, 2018.
% verified: publisher URL https://www.amazon.com/dp/B0HKPL1X1S 2026-09-27 (Independently published, 22 September 2026, 488 pages, ISBN 9798176411898; title and subtitle from the title page of the interior PDF, whose copyright page reads "First edition, 2026" and prints no publisher or ISBN; pinpoints follow this edition, not the draft of about 195 pages)
\bibitem{P26} N. G. Polson. \emph{On Hilbert's 8th Problem: the reciprocal Xi clock, Thorin measures, van Dantzig pairs, zero free approximations, hitting times and the reciprocal zeta function}, first edition. Independently published, 2026.
% verified: manuscript N/A 2026-09-27 (title, sole author and date from the source file eta_paper.tex; no arXiv record)
\bibitem{P26eta} N. G. Polson. Thorin measures of the Dedekind eta function. Manuscript, September 2026.
% verified: manuscript N/A 2026-10-14 (title from the source file chapter_max_wald_thorin.tex, received on 14 October 2026 as chapter_max_wald.zip; sole author as for the first author's other manuscripts; no arXiv record)
\bibitem{P26mw} N. G. Polson. The P\'olya density as a maximum order statistic and the Thorin measure of the reciprocal Xi function. Manuscript, October 2026.
% verified: manuscript N/A 2026-09-27 (title, byline and date from reciprocal-zeta/paper/reciprocal_zeta_thorin.tex)
\bibitem{P26rz} N. G. Polson. Thorin type representations of the reciprocal zeta function on the critical line. Manuscript, September 2026, now Part IX of the first edition of \emph{On Hilbert's 8th Problem} (2026).
% verified: manuscript N/A 2026-09-27 (authors' source thorin_xi_computation.tex, title page dated September 2026; not on arXiv)
\bibitem{PS26} N. G. Polson and V. Sokolov. Computing the Thorin measure of the reciprocal Xi function: a P\'olya route and a total variation comparison with the Riemann atoms. Manuscript, September 2026.
% verified: manuscript N/A 2026-10-14 (working report of both authors, dated 7 October 2026, cited by the source files of Chapter 14 received on 14 October 2026 as crh_ch14.zip, which supersede it; no arXiv record)
\bibitem{FA26} N. G. Polson and V. Sokolov. Four attacks on the computational path: a working report. Manuscript, 7 October 2026.
% verified: CrossRef 10.1007/bf02565336 2026-09-26
\bibitem{Polya} G. P\'olya. Bemerkung \"uber die Integraldarstellung der Riemannschen $\xi$-Funktion. \emph{Acta Mathematica} 48 (1926), 305--317.
% verified: CrossRef 10.1007/s40687-019-0193-1 2026-09-26
\bibitem{Polymath} D. H. J. Polymath. Effective approximation of heat flow evolution of the Riemann $\xi$ function, and a new upper bound for the de Bruijn--Newman constant. \emph{Research in the Mathematical Sciences} 6 (2019), paper 31.
% verified: CrossRef 10.1007/bf02418544 2026-09-26
\bibitem{Riesz16} M. Riesz. Sur l'hypoth\`ese de Riemann. \emph{Acta Mathematica} 40 (1916), 185--190.
% verified: CrossRef 10.2307/2308012 2026-09-28; pages 26 to 29 from MathSciNet MR Lookup, 2026-10-02
\bibitem{Robbins55} H. Robbins. A remark on Stirling's formula. \emph{American Mathematical Monthly} 62 (1955), 26--29.
% verified: CrossRef 10.1017/fmp.2020.6 2026-09-26
\bibitem{RodgersTao} B. Rodgers and T. Tao. The de Bruijn--Newman constant is non-negative. \emph{Forum of Mathematics, Pi} 8 (2020), e6.
% verified: CrossRef 10.1215/ijm/1255631807 2026-09-28 (pages 64 to 94 from zbMATH 0122.05001; CrossRef gives no pages)
\bibitem{RS62} J. B. Rosser and L. Schoenfeld. Approximate formulas for some functions of prime numbers. \emph{Illinois Journal of Mathematics} 6 (1962), 64--94.
% verified: zbMATH 0792.11034 2026-09-26; pages 415 to 440 from zbMATH 0792.11034, which adds the supplement S17 to S23, 2026-10-02
\bibitem{Rumely93} R. Rumely. Numerical computations concerning the ERH. \emph{Mathematics of Computation} 61 (1993), 415--440.
% verified: zbMATH 0050.08003 2026-09-26
\bibitem{Salem53} R. Salem. Sur une proposition \'equivalente \`a l'hypoth\`ese de Riemann. \emph{Comptes Rendus de l'Acad\'emie des Sciences} 236 (1953), 1127--1128.
% verified: zbMATH 0973.60001 2026-09-26
\bibitem{Sato99} K. Sato. \emph{L\'evy Processes and Infinitely Divisible Distributions}. Cambridge Studies in Advanced Mathematics 68, Cambridge University Press, 1999.
% verified: CrossRef 10.1515/9783110269338 2026-09-26
\bibitem{boSSV} R. L. Schilling, R. Song and Z. Vondra\v{c}ek. \emph{Bernstein Functions: Theory and Applications}, second edition. Walter de Gruyter, 2012.
% verified: CrossRef 10.1007/bf02790092 2026-09-27 (accent in Polya from zbMATH 0045.37602 and the original first page; CrossRef omits it)
\bibitem{Schoenberg51} I. J. Schoenberg. On P\'olya frequency functions. I. The totally positive functions and their Laplace transforms. \emph{Journal d'Analyse Math\'ematique} 1 (1951), 331--374.
% verified: zbMATH 0061.08402 2026-09-26
\bibitem{Selberg46} A. Selberg. Contributions to the theory of the Riemann zeta-function. \emph{Archiv for Mathematik og Naturvidenskab} 48 (1946), 89--155.
% verified: CrossRef 10.1007/s10998-010-1037-3 2026-09-26
\bibitem{SD10} J. Sondow and C. Dumitrescu. A monotonicity property of Riemann's xi function and a reformulation of the Riemann hypothesis. \emph{Periodica Mathematica Hungarica} 60 (2010), 37--40.
% verified: CrossRef 10.1090/s0025-5718-1994-1254148-8 2026-10-11 (Math. Comp. 63, no. 208, first page 747; CrossRef gives no last page); 747--748 as in Ivic, arXiv:math/0311162, ref. [58]
\bibitem{Spira94} R. Spira. Some zeros of the Titchmarsh counterexample. \emph{Mathematics of Computation} 63 (1994), 747--748.
% verified: CrossRef 10.1214/aoms/1177698806 2026-10-09
\bibitem{Steutel67} F. W. Steutel. Note on the infinite divisibility of exponential mixtures. \emph{Annals of Mathematical Statistics} 38 (1967), 1303--1305.
% verified: CrossRef 10.1093/imamat/23.1.97 2026-09-26
\bibitem{Talbot79} A. Talbot. The accurate numerical inversion of Laplace transforms. \emph{Journal of the Institute of Mathematics and its Applications} 23 (1979), 97--120.
% verified: zbMATH 1492.11154 2026-09-27 (title as printed on p. 305; zbMATH adds a comma after fields)
\bibitem{Tate67} J. T. Tate. Fourier analysis in number fields and Hecke's zeta-functions. In \emph{Algebraic Number Theory}, J. W. S. Cassels and A. Fr\"ohlich, eds., Academic Press, 1967, 305--347.
% verified: zbMATH 0601.10026 2026-09-26
\bibitem{Titchmarsh} E. C. Titchmarsh. \emph{The Theory of the Riemann Zeta-Function}, second edition, revised by D. R. Heath-Brown. Oxford University Press, 1986.
% verified: CrossRef 10.1016/j.jnt.2013.07.017 2026-09-26
\bibitem{Trudgian} T. S. Trudgian. An improved upper bound for the argument of the Riemann zeta-function on the critical line II. \emph{Journal of Number Theory} 134 (2014), 280--292.
% verified: 2 November 2026 at arxiv.org/abs/2206.03682 (abstract read; the sign convention of the screw function is not in the abstract)
\bibitem{Suzuki22} M. Suzuki. Aspects of the screw function corresponding to the Riemann zeta function. arXiv:2206.03682.
% verified: zbMATH 0049.03205 2026-09-26; pages 252 to 265 from zbMATH 0049.03205, 2026-10-02
\bibitem{Weil52} A. Weil. Sur les ``formules explicites'' de la th\'eorie des nombres premiers. \emph{Communications du S\'eminaire Math\'ematique de l'Universit\'e de Lund}, tome suppl\'ementaire (1952), 252--265.
% verified: zbMATH 0063.08245 2026-09-27
\bibitem{Widder41} D. V. Widder. \emph{The Laplace Transform}. Princeton Mathematical Series 6, Princeton University Press, 1941.
% verified: CrossRef 10.2307/1968102 2026-09-26; pages 1 to 100 from zbMATH 0004.05905, 2026-10-02
\bibitem{Wiener32} N. Wiener. Tauberian theorems. \emph{Annals of Mathematics} 33 (1932), 1--100.
% verified: CrossRef 10.1215/s0012-7094-56-02317-1 2026-09-26; pages 189 to 207 from zbMATH 0070.28501, 2026-10-02
\bibitem{Williamson} R. E. Williamson. Multiply monotone functions and their Laplace transforms. \emph{Duke Mathematical Journal} 23 (1956), 189--207.
% verified: CrossRef 10.1214/aop/1176995474 2026-09-26; pages 523 to 531 from zbMATH 0394.60017, 2026-10-02
\bibitem{Yamazato78} M. Yamazato. Unimodality of infinitely divisible distribution functions of class $L$. \emph{Annals of Probability} 6 (1978), 523--531.
% verified: zbMATH 0817.11041 2026-09-27 (DOI 10.2969/aspm/02110281; editors from zbMATH 0771.00036; conference at Tokyo Institute of Technology, August 14 to 18, 1990, per Project Euclid)
\bibitem{Yoshida92} H. Yoshida. On Hermitian forms attached to zeta functions. In \emph{Zeta Functions in Geometry (Tokyo, 1990)}, N. Kurokawa and T. Sunada, eds., Advanced Studies in Pure Mathematics 21, Kinokuniya, 1992, 281--325.
% verified: CrossRef 10.1007/s00208-006-0051-5 2026-11-02
\bibitem{Zhang07} W.-B. Zhang. Beurling primes with RH and Beurling primes with large oscillation. \emph{Mathematische Annalen} 337 (2007), 671--704.
% verified: arXiv 2608.24827 2026-10-17; author Zhu, Xuefeng and title match v2 (2026-09-02) via the arXiv API and zbMATH; v1 (2026-08-25) lists Chuk, Marcus alone with a different title, so the author of record changed between versions (the arXiv comment of v2 says: v1 was a 9 page announcement, author name and affiliation updated); the book cites v2, the full version of 34 pages
\bibitem{Zhu26} X. Zhu. Weil positivity in compact windows: a finite reduction, certified two-sided bounds, and a Landau--Widom decay law. arXiv:2608.24827v2, 2 September 2026.
\bibitem{Weil48} A. Weil. \emph{Sur les courbes alg\'ebriques et les vari\'et\'es qui s'en d\'eduisent}. Actualit\'es scientifiques et industrielles 1041, Hermann, Paris, 1948. % verified: standard reference for Weil's proof of RH for curves; book not reread for this chapter
% verified: CrossRef 10.1090/s0002-9947-1957-0091566-1 2026-11-16
\bibitem{KarlinMcGregor57} S. Karlin and J. McGregor. The differential equations of birth-and-death processes, and the Stieltjes moment problem. \emph{Transactions of the American Mathematical Society} 85 (1957), 489--546.
% verified: CrossRef 10.1090/s0002-9939-99-05200-4 2026-11-16 (WARN: CrossRef issued date is 1999, volume 128 is 2000; draft year 2000 kept)
\bibitem{Gneiting00} T. Gneiting. Kuttner's problem and a P\'olya type criterion for characteristic functions. \emph{Proceedings of the American Mathematical Society} 128 (2000), 1721--1728.

% verified: CrossRef 10.1090/s0002-9939-01-05839-7 2026-11-16
\bibitem{Gneiting01} T. Gneiting. Criteria of P\'olya type for radial positive definite functions. \emph{Proceedings of the American Mathematical Society} 129 (2001), 2309--2318.

% verified: CrossRef 10.1016/j.spl.2023.109889 2026-11-16 (WARN: the draft cites arXiv:2306.08317, which the arXiv record confirms with the same title and authors; a journal version now exists, so the journal form is written here)
\bibitem{NakamuraSuzuki23} T. Nakamura and M. Suzuki. On infinitely divisible distributions related to the Riemann hypothesis. \emph{Statistics \& Probability Letters} 201 (2023), 109889.

% UNCHECKED: second pass: no CrossRef record; OpenAlex gives a 1949 Project Euclid record only, author and pages not retrieved
\bibitem{Polya49} G. P\'olya. Remarks on characteristic functions. \emph{Proceedings of the First Berkeley Symposium on Mathematical Statistics and Probability} (1949), 115--123.

% UNCHECKED: title and journal from the draft, pages 864--867 from the referee note NAK2 (not looked up)
\bibitem{Getoor79} R. K. Getoor. The Brownian escape process. \emph{Annals of Probability} 7 (1979), 864--867.
% New bibliography entries of ch/nc_v3_gamma.tex (Volume III, Chapter 5), to be merged into the Volume III bibliography.
% verified: CrossRef 10.1007/bf01231885 2026-11-16
\bibitem{Deninger92} C. Deninger. Local $L$-factors of motives and regularized determinants. \emph{Inventiones Mathematicae} 107 (1992), 135--150.

% verified: CrossRef 10.1017/cbo9780511526251 2026-11-16 (WARN: the record gives Cambridge University Press 2004 with no edition field; "second edition" from the draft is not confirmed and is omitted)
\bibitem{GasperRahman04} G. Gasper and M. Rahman. \emph{Basic Hypergeometric Series}. Cambridge University Press, 2004.

% UNCHECKED: second pass: CrossRef 1904 Jackson hits are other papers (Neumann expansion; basic gamma 1905); the 1904 gamma paper not found
\bibitem{Jackson04} F. H. Jackson. A generalization of the functions $\Gamma(n)$ and $x^n$. \emph{Proceedings of the Royal Society of London} 74 (1904), 64--72.


% verified: CrossRef 10.1016/j.spl.2023.109889 2026-11-16 (WARN: the draft cites arXiv:2306.08317, which the arXiv record confirms with the same title and authors; a journal version now exists, so the journal form is written here)
% Works named in Volume III prose (citation-verifier, 2026-11-17)
% New bibliography entries for Volume III. Keys checked against ch/bib.tex and all ch/nc_*_bib.tex: no clashes.
% Already in the book under other keys: GM26 (Guth and Maynard), Montgomery73, RS62 (Rosser and Schoenfeld 1962), JessenWintner35 (the 1935 Wintner paper is joint with Jessen), Gantmacher59 and Polya49 (in the other Volume III bib files, still unchecked).
% Still without an entry: Goldfeld (1982), Diamond and Erdos, Hardy-Littlewood-Polya book (HLP34 exists in nc_v3_classics_bib.tex, unchecked).
% verified: CrossRef 10.5802/jtnb.22 2026-11-17
\bibitem{Odlyzko90} A. M. Odlyzko. Bounds for discriminants and related estimates for class numbers, regulators and zeros of zeta functions: a survey of recent results. \emph{Journal de Th\'eorie des Nombres de Bordeaux} 2 (1990), 119--141.

% verified: CrossRef 10.1007/bf02592003 2026-11-17
\bibitem{Bohman61} H. Bohman. Approximate Fourier analysis of distribution functions. \emph{Arkiv f\"or Matematik} 4 (1961), 99--157.

% UNCHECKED: CrossRef gives the title, author and DOI 10.1007/bfb0096066 (Lecture Notes in Mathematics, pages 136--153), but no year, and the venue is a Bourbaki seminar, not the Delange-Pisot-Poitou seminar the prose names; confirm venue and year before use
\bibitem{Poitou77} G. Poitou. Minorations de discriminants [d'apr\`es A. M. Odlyzko]. In \emph{S\'eminaire Bourbaki}, Lecture Notes in Mathematics, pages 136--153. Springer. \url{https://doi.org/10.1007/bfb0096066}

% verified: CrossRef 10.4153/cjm-2005-014-0 2026-11-17
\bibitem{KuoMurty05} W. Kuo and M. R. Murty. On a conjecture of Birch and Swinnerton-Dyer. \emph{Canadian Journal of Mathematics} 57 (2005), 328--337.

% verified: CrossRef 10.4064/aa-91-3-209-228 2026-11-17 (the 1993 Bull. Austral. Math. Soc. paper "Small zeros of quadratic L-functions" is a different item; confirm the prose means this one)
\bibitem{OzlukSnyder99} A. E. \"Ozl\"uk and C. Snyder. On the distribution of the nontrivial zeros of quadratic $L$-functions close to the real axis. \emph{Acta Arithmetica} 91 (1999), 209--228.

% verified: CrossRef 10.2307/2661390 2026-11-17
\bibitem{Soundararajan00} K. Soundararajan. Nonvanishing of quadratic Dirichlet $L$-functions at $s=\frac12$. \emph{Annals of Mathematics} 152 (2000), 447--488.

% verified: CrossRef 10.1007/bf01388809 2026-11-17
\bibitem{GrossZagier86} B. H. Gross and D. B. Zagier. Heegner points and derivatives of $L$-series. \emph{Inventiones Mathematicae} 84 (1986), 225--320.

% verified: CrossRef 10.1080/10586458.1994.10504289 2026-11-17
\bibitem{RubinsteinSarnak94} M. Rubinstein and P. Sarnak. Chebyshev's bias. \emph{Experimental Mathematics} 3 (1994), 173--197.

% verified: CrossRef 10.1016/0041-5553(69)90124-4 2026-11-17
\bibitem{Uvarov69} V. B. Uvarov. The connection between systems of polynomials orthogonal with respect to different distribution functions. \emph{USSR Computational Mathematics and Mathematical Physics} 9 (1969), 25--36.
% New bibliography entries for ch/nc_v3_classics.tex. Taken from the draft's list (misc/reviews/nick-ch17-classics/chapter17/chH_classics.tex); none of these appears in misc/vol3/bib_new.tex except Lehmer88, so each is UNCHECKED against an external source.
% UNCHECKED: from the draft reference list
\bibitem{vonMangoldt95} H. von Mangoldt, Zu Riemanns Abhandlung \"Uber die Anzahl der Primzahlen unter einer gegebenen Gr\"osse, \emph{J. Reine Angew. Math.} 114 (1895), 255 to 305.
% verified: OpenAlex 10.1112/plms/s2_14.1.347 2026-11-17 (record: Proc. London Math. Soc. (2) 14 (1915), first page 347; pages to 409 from the draft)
\bibitem{Ramanujan15} S. Ramanujan, Highly composite numbers, \emph{Proc. London Math. Soc.} (2) 14 (1915), 347 to 409.
% verified: CrossRef 10.1023/a:1009764017495 2026-11-17
\bibitem{Ramanujan97} S. Ramanujan, Highly composite numbers, annotated and with a foreword by J. L. Nicolas and G. Robin, \emph{Ramanujan J.} 1 (1997), 119 to 153.
% verified: arXiv 1211.6944 2026-11-17 (WARN: arXiv authors are Nicolas and Sondow, the draft lists Nicolas alone; published as Contemp. Math. 627 (2014), CrossRef 10.1090/conm/627/12539)
\bibitem{Nicolas12} J. L. Nicolas and J. Sondow, Ramanujan, Robin, highly composite numbers, and the Riemann hypothesis, \emph{Contemp. Math.} 627 (2014); arXiv:1211.6944.
% UNCHECKED: from the draft reference list
\bibitem{Robin84} G. Robin, Grandes valeurs de la fonction somme des diviseurs et hypoth\`ese de Riemann, \emph{J. Math. Pures Appl.} 63 (1984), 187 to 213.
% UNCHECKED: from the draft reference list
\bibitem{Littlewood14} J. E. Littlewood, Sur la distribution des nombres premiers, \emph{C. R. Acad. Sci. Paris} 158 (1914), 1869 to 1872.
% UNCHECKED: from the draft reference list
\bibitem{HLP34} G. H. Hardy, J. E. Littlewood and G. P\'olya, \emph{Inequalities}, Cambridge University Press, 1934; second edition 1952.
% verified: CrossRef 10.1112/jlms/s2-8.1.73 2026-11-17
\bibitem{MV74} H. L. Montgomery and R. C. Vaughan, Hilbert's inequality, \emph{J. London Math. Soc.} (2) 8 (1974), 73 to 82.
% verified: CrossRef 10.1090/s0025-5718-1988-0917834-x 2026-11-16
\bibitem{Lehmer88} D. H. Lehmer. The sum of like powers of the zeros of the Riemann zeta function. \emph{Mathematics of Computation} 50 (1988), 265--273.
% Bibliography entries for ch/nc_v3_cheb.tex that are not in ch/bib.tex. 17 Nov 2026.
% UNCHECKED: not in misc/vol3/bib_new.tex; entered from the standard citation (Titchmarsh, Section 14.25 cites it); verify volume and pages before release.
\bibitem{Littlewood12} J. E. Littlewood. Quelques cons\'equences de l'hypoth\`ese que la fonction $\zeta(s)$ de Riemann n'a pas de z\'eros dans le demi-plan $\Re s>\frac12$. \emph{Comptes Rendus de l'Acad\'emie des Sciences Paris} 154 (1912), 263--266.

% New bibliography entry of Chapter nc_v3_cw (prefix cw:), to be pooled with ch/bib.tex by the main session. Not found in ch/bib.tex, ch/*_bib.tex or misc/vol3/bib_new.tex.
\bibitem{DunlopNewman75} F. Dunlop and C. M. Newman. Multicomponent field theories and classical rotators. \emph{Communications in Mathematical Physics} 44 (1975), 223--235. % verified: CrossRef, DOI 10.1007/bf01609827
% verified: CrossRef 10.1515/crll.1989.399.1 2026-11-16
\bibitem{Conrey89} J. B. Conrey. More than two fifths of the zeros of the Riemann zeta function are on the critical line. \emph{J. reine angew. Math.} 399 (1989), 1--26.
% verified against the draft note only (fields as in misc/reviews/nick-route-to-rh/route/route_chapter.tex); pages not given
% verified: OpenAlex title search "Sulla determinazione empirica di una legge di distribuzione" (volume 4, pages 83-91, 1933; no DOI) 2026-11-17
\bibitem{Kolmogorov33} A. Kolmogorov. Sulla determinazione empirica di una legge di distribuzione. \emph{Giornale dell'Istituto Italiano degli Attuari} 4 (1933), 83--91.
% UNCHECKED
\bibitem{Prekopa73} A. Pr\'ekopa. On logarithmic concave measures and functions. \emph{Acta Scientiarum Mathematicarum (Szeged)} 34 (1973).
\bibitem{NCIwaniecConv06} H. Iwaniec. Conversations on the exceptional character. In \emph{Analytic Number Theory} (Cetraro, 2002), Lecture Notes in Mathematics 1891, Springer, 2006, 97--132. % verified: DOI 10.1007/978-3-540-36364-4_3, as supplied by the authors; the text was not read, and we do not know whether it describes the elliptic curve family of Section~1
\bibitem{NCConreyIwaniec02} J. B. Conrey and H. Iwaniec. Spacing of zeros of Hecke $L$-functions and the class number problem. \emph{Acta Arithmetica} 103 (2002), 259--312. % verified: DOI 10.4064/aa103-3-5, as supplied by the authors
\bibitem{NCConrey03} J. B. Conrey. The Riemann hypothesis. \emph{Notices of the American Mathematical Society} 50 (2003), 341--353. % UNCHECKED: fields from Conrey's publication list only; the survey text was not read
\bibitem{NCMestre86} J.-F. Mestre. Formules explicites et minorations de conducteurs de vari\'et\'es alg\'ebriques. \emph{Compositio Mathematica} 58 (1986), 209--232. % verified: Numdam CM_1986__58_2_209_0 and EuDML 89769, 2026-11-04 (no DOI)
\bibitem{NCBrumer92} A. Brumer. The average rank of elliptic curves I. \emph{Inventiones Mathematicae} 109 (1992), 445--472. % verified: CrossRef 10.1007/bf01232033 2026-11-04
\bibitem{NCBoas54} R. P. Boas. \emph{Entire Functions}. Academic Press, New York, 1954. % verified: publisher and year, WorldCat 847696, 2026-11-04; pages 124 to 132 for Akhiezer's factorisation per secondary sources, book not read
\bibitem{NCGrossZagier86} B. H. Gross and D. B. Zagier. Heegner points and derivatives of $L$-series. \emph{Inventiones Mathematicae} 84 (1986), 225--320. % verified: CrossRef 10.1007/BF01388809 2026-11-04
\bibitem{NCKolyvagin89} V. A. Kolyvagin. Finiteness of $E(\mathbb Q)$ and $\mathrm{Sha}(E,\mathbb Q)$ for a subclass of Weil curves. \emph{Mathematics of the USSR Izvestiya} 32 (1989), 523--541. % verified: CrossRef 10.1070/im1989v032n03abeh000779 2026-11-04 (English translation of the 1988 Izvestiya paper)
\bibitem{BohrCourant} H. Bohr and R. Courant. Neue Anwendungen der Theorie der Diophantischen Approximationen auf die Riemannsche Zetafunktion. \emph{Journal f\"ur die reine und angewandte Mathematik} 144 (1914), 249--274. % UNCHECKED: fields from the bibliography of the note, not seen in the original
\bibitem{NCOAI26c} OpenAI. Uniform exclusion of Landau--Siegel zeros. Preprint dated 1 October 2026, \url{https://github.com/openai/math} (family 003). % verified: CONTENTS.md and lean/docs/003.md of the repository, read 28 October 2026
\bibitem{NCQuanta26} Quanta Magazine. AI proved the quasi-Riemann hypothesis. What happens now? 19 October 2026, \url{https://www.quantamagazine.org/ai-proved-the-quasi-riemann-hypothesis-what-happens-now-20261019}. % verified: article fetched 28 October 2026
\bibitem{NCDunnRadziwill24} A. Dunn and M. Radziwi\l\l. Bias in cubic Gauss sums: Patterson's conjecture. \emph{Annals of Mathematics} (2) 200 (2024), 967--1057. % verified: journal page annals.math.princeton.edu/2024/200-3/p03, DOI 10.4007/annals.2024.200.3.3
\bibitem{NCPatterson77} S. J. Patterson. A cubic analogue of the theta series. \emph{Journal f\"ur die reine und angewandte Mathematik} 296 (1977), 125--161. % verified: reference [38] of OAI26b with DOI 10.1515/crll.1977.296.125; journal page not opened
\bibitem{NCKubota69} T. Kubota. \emph{On Automorphic Functions and the Reciprocity Law in a Number Field}. Lectures in Mathematics, Department of Mathematics, Kyoto University, No.~2, Kinokuniya, Tokyo, 1969. % verified: reference [26] of OAI26b only; UNCHECKED against a library record
\bibitem{NCGoldmakherLouvel13} L. Goldmakher and B. Louvel. A quadratic large sieve inequality over number fields. \emph{Mathematical Proceedings of the Cambridge Philosophical Society} 154 (2013), 193--212. % verified: reference [13] of OAI26b with DOI 10.1017/S0305004112000370; journal page not opened
% verified: CrossRef 10.1137/1.9781611975949 2026-11-16 (WARN: CrossRef lists the Extended Edition, SIAM 2019; the draft's 2013 date is not in CrossRef, so 2019 is written here)
\bibitem{Trefethen19} L. N. Trefethen. \emph{Approximation Theory and Approximation Practice}, Extended Edition. SIAM, 2019.
% verified: CrossRef 10.1090/trans2/103/02 2026-11-16
\bibitem{KacKrein74} I. S. Kac and M. G. Krein. On the spectral functions of the string. \emph{American Mathematical Society Translations} (2) 103 (1974), 19--102.

% verified: CrossRef 10.5802/aif.2125 2026-11-16
\bibitem{RoynetteYor05} B. Roynette and M. Yor. Couples de Wald ind\'efiniment divisibles. Exemples li\'es \`a la fonction gamma d'Euler et \`a la fonction zeta de Riemann. \emph{Annales de l'Institut Fourier} 55 (2005), 1219--1283.
% verified: CrossRef 10.1090/s0002-9947-1957-0091566-1 2026-11-16
% verified: CrossRef 10.5802/afst.108 2026-11-16 (CrossRef gives pages J1--J122, the draft gives 1 to 122); also entered by the Chapter 13 editor if that chapter cites it, remove one copy
\bibitem{Stieltjes94} T. J. Stieltjes. Recherches sur les fractions continues. \emph{Annales de la Facult\'e des Sciences de Toulouse} 8 (1894), J1--J122.

% verified: CrossRef 10.4064/aa140-4-4 2026-11-16 (WARN: the draft's arXiv:1810.10426 is a different paper, Zaghloul; this journal record matches the draft title and authors)
\bibitem{SaiasWeingartner09} E. Saias and A. Weingartner. Zeros of Dirichlet series with periodic coefficients. \emph{Acta Arithmetica} 140 (2009), 335--344.

% UNCHECKED: second pass: CrossRef lists only reviews of the book; Chelsea 1959 record not found
\bibitem{Gantmacher59} F. R. Gantmacher. \emph{The Theory of Matrices}, Vol. 2. Chelsea, 1959, Chapter XV.
% UNCHECKED: second pass: CrossRef and OpenAlex hits are unrelated Ramanujan papers; the 1915 note is not indexed
\bibitem{Ram15} S. Ramanujan, New expressions for Riemann's functions $\xi(s)$ and $\Xi(t)$, \emph{Quart. J. Math.} 46 (1915), 253 to 260.
% UNCHECKED: second pass: CrossRef and OpenAlex hits are unrelated Ramanujan papers; the 1916 note is not indexed
\bibitem{RamF} S. Ramanujan, Some formulae in the analytic theory of numbers, \emph{Messenger Math.} 45 (1916), 81 to 84.
% UNCHECKED: second pass: CrossRef and OpenAlex hits are unrelated Ramanujan papers; the 1915 note is not indexed
\bibitem{RamG} S. Ramanujan, Some definite integrals connected with Gauss's sums, \emph{Messenger Math.} 44 (1915), 75 to 85.
% verified (misc/vol3/bib_new.tex)
\bibitem{Keiper96} J. B. Keiper. On the zeros of the Ramanujan $\tau$-Dirichlet series in the critical strip. \emph{Mathematics of Computation} 65 (1996), 1613--1619.

% New bibliography entries of Chapter "Barndorff-Nielsen's tools on the ladder" (ch/nc_v3_bn.tex, prefix bn:), to be merged into ch/bib.tex by the main session. Existing keys used: BNKS82, BNMS06, BNT02, BP99, BPY01, CNV86, Kristiansen94, Lagarias, Polya, Montgomery73, Dyson62. None of the entries below was found in misc/vol3/bib_new.tex or ch/bib.tex. The fields are those of the draft (misc/reviews/nick-bn/bn/bn_chapter.tex); no external check was made in this pass, and the draft README reports a web check only for the 1977, 1978, 1982, 2002, 2004, 2006 Barndorff-Nielsen papers, Benaych-Georges, Perez-Abreu and Sakuma, and Berg, Koumandos and Pedersen.
% verified: CrossRef 10.1098/rspa.1977.0041 2026-11-18
\bibitem{BN77} O. Barndorff-Nielsen. Exponentially decreasing distributions for the logarithm of particle size. \emph{Proceedings of the Royal Society of London A} 353 (1977), 401--419.

% verified: CrossRef 10.1007/bf00533162 2026-11-18
\bibitem{BNH77} O. Barndorff-Nielsen and C. Halgreen. Infinite divisibility of the hyperbolic and generalized inverse Gaussian distributions. \emph{Zeitschrift f\"ur Wahrscheinlichkeitstheorie und verwandte Gebiete} 38 (1977), 309--312.

% verified: CrossRef 10.1016/0304-4149(78)90036-4 2026-11-18
\bibitem{BBH78} O. Barndorff-Nielsen, P. Bl{\ae}sild and C. Halgreen. First hitting time models for the generalized inverse Gaussian distribution. \emph{Stochastic Processes and their Applications} 7 (1978), 49--54.

% verified: CrossRef 10.1111/1467-9868.00282 2026-11-18
\bibitem{BNS01} O. E. Barndorff-Nielsen and N. Shephard. Non-Gaussian Ornstein-Uhlenbeck-based models and some of their uses in financial economics. \emph{Journal of the Royal Statistical Society B} 63 (2001), 167--241.

% UNCHECKED: CrossRef 10.4213/tvp464 gives the Russian original, year 2000 (draft says 2001, English translation); year and journal not changed
\bibitem{BN01} O. E. Barndorff-Nielsen. Superposition of Ornstein-Uhlenbeck type processes. \emph{Theory of Probability and its Applications} 45 (2001).

% verified: CrossRef 10.1016/j.spa.2005.09.011 2026-11-18
\bibitem{BNT06} O. E. Barndorff-Nielsen and S. Thorbj{\o}rnsen. Regularizing mappings of L\'evy measures. \emph{Stochastic Processes and their Applications} 116 (2006), 423--446.

% verified: CrossRef 10.1214/009117904000000982 2026-11-18
\bibitem{BenaychGeorges05} F. Benaych-Georges. Classical and free infinitely divisible distributions and random matrices. \emph{Annals of Probability} 33 (2005), 1134--1170.

% verified: CrossRef 10.1214/ejp.v10-246 2026-11-18
\bibitem{CabanalDuvillard05} T. Cabanal-Duvillard. A matrix representation of the Bercovici-Pata bijection. \emph{Electronic Journal of Probability} 10 (2005), 632--661.

% UNCHECKED: CrossRef 10.1512/iumj.2014.63.5288 lists Thorbjornsen first, Haagerup second (draft has Haagerup first); author order not changed
\bibitem{HaagerupThorbjornsen14} U. Haagerup and S. Thorbj{\o}rnsen. On the free gamma distributions. \emph{Indiana University Mathematics Journal} 63 (2014).

% UNCHECKED: arXiv:1905.04131 found (title matches, 2019); journal version exists, Math. Nachr. 294 (2020) 426--449, DOI 10.1002/mana.201900217; not changed
\bibitem{BKP19} C. Berg, S. Koumandos and H. L. Pedersen. Nielsen's beta function and some infinitely divisible distributions. arXiv:1905.04131, 2019.

% UNCHECKED: found in CrossRef 10.1215/ijm/1256055197 (year, volume, title match) but CrossRef gives no pages; pages 9--23 unconfirmed
\bibitem{Grosswald66} E. Grosswald. Generalization of a formula of Hayman and its application to the study of Riemann's zeta function. \emph{Illinois Journal of Mathematics} 10 (1966), 9--23.

% verified: CrossRef 10.1515/crll.1956.196.67 2026-11-18
\bibitem{Hayman56} W. K. Hayman. A generalisation of Stirling's formula. \emph{Journal f\"ur die reine und angewandte Mathematik} 196 (1956), 67--95.

% verified: CrossRef 10.1007/bf00533993 2026-11-18
\bibitem{Grosswald76a} E. Grosswald. The Student $t$-distribution of any degree of freedom is infinitely divisible. \emph{Zeitschrift f\"ur Wahrscheinlichkeitstheorie und verwandte Gebiete} 36 (1976), 103--109.

% verified: CrossRef 10.1214/aop/1176996038 2026-11-18
\bibitem{Grosswald76b} E. Grosswald. The Student $t$-distribution for odd degrees of freedom is infinitely divisible. \emph{Annals of Probability} 4 (1976), 680--683.

% verified: CrossRef 10.1090/s0002-9947-1967-0202685-7 2026-11-18
\bibitem{Grosswald67} E. Grosswald. Oscillation theorems of arithmetical functions. \emph{Transactions of the American Mathematical Society} 126 (1967), 1--28.

% verified: CrossRef 10.2140/pjm.1978.74.357 2026-11-18
\bibitem{GrosswaldSchnitzer78} E. Grosswald and F. J. Schnitzer. A class of modified $\zeta$ and $L$ functions. \emph{Pacific Journal of Mathematics} 74 (1978), 357--364.

% verified: CrossRef 10.4064/aa-9-4-365-373 2026-11-18
\bibitem{BatemanGrosswald64} P. T. Bateman and E. Grosswald. On Epstein's zeta function. \emph{Acta Arithmetica} 9 (1964), 365--373.

% verified: CrossRef 10.1007/bf00533246 2026-11-18
\bibitem{Halgreen79} C. Halgreen. Self-decomposability of the generalized inverse Gaussian and hyperbolic distributions. \emph{Zeitschrift f\"ur Wahrscheinlichkeitstheorie und verwandte Gebiete} 47 (1979), 13--17.

% verified: CrossRef 10.1214/aop/1176995766 2026-11-18
\bibitem{Ismail77} M. E. H. Ismail. Bessel functions and the infinite divisibility of the Student $t$-distribution. \emph{Annals of Probability} 5 (1977), 582--585.
\bibitem{NCHamburger21} H. Hamburger. \"Uber die Riemannsche Funktionalgleichung der $\zeta$-Funktion. \emph{Mathematische Zeitschrift} 10 (1921), 240--254. % verified: 10.1007/bf01211612 2026-11-09
\bibitem{NCBeurling37} A. Beurling. Analyse de la loi asymptotique de la distribution des nombres premiers g\'en\'eralis\'es. I. \emph{Acta Mathematica} 68 (1937), 255--291. % verified: 10.1007/bf02546666 2026-11-09
\bibitem{NCDMV06} H. G. Diamond, H. L. Montgomery and U. M. A. Vorhauer. Beurling primes with large oscillation. \emph{Mathematische Annalen} 334 (2006), 1--36. % verified: 10.1007/s00208-005-0638-2 2026-11-09
\bibitem{NCDiamondZhang16} H. G. Diamond and W.-B. Zhang. \emph{Beurling Generalized Numbers}. Mathematical Surveys and Monographs 213, American Mathematical Society, 2016. % verified: 10.1090/surv/213 2026-11-09
\bibitem{NCCohnElkies03} H. Cohn and N. Elkies. New upper bounds on sphere packings I. \emph{Annals of Mathematics} 157 (2003), 689--714. % verified: 10.4007/annals.2003.157.689 2026-11-09
\bibitem{NCKP99} J. Kaczorowski and A. Perelli. On the structure of the Selberg class, I: $0\le d\le1$. \emph{Acta Mathematica} 182 (1999), 207--241. % verified: 10.1007/bf02392574 2026-11-09
\bibitem{NCGrenanderSzego58} U. Grenander and G. Szeg\H o. \emph{Toeplitz Forms and Their Applications}. University of California Press, Berkeley, 1958. % verified: CrossRef 10.1525/9780520355408 2026-11-10
\end{thebibliography}
